\begin{itemize}
  \item \textbf{Assume students are new to standard infinite set symbols.} Do not assume prior familiarity with $\mathbb{N}$, $\mathbb{Z}$, $\mathbb{Q}$, $\mathbb{R}$, or $\mathbb{C}$.

  \item \textbf{You may use $\mathbb{W}$ and/or $\mathbb{N}$, but define them first.} If using these in examples or exercises, explicitly define them before use, e.g.
  $\mathbb{W} = \{0,1,2,3,\ldots\}$ (whole numbers) and
  $\mathbb{N} = \{1,2,3,\ldots\}$ (natural numbers).

  \item \textbf{Prefer multiple-choice when symbols are required.} Since symbols like $\in$, $\notin$, and $\mathbb{W}$/$\mathbb{N}$ are difficult to type, present questions involving these symbols in multiple-choice form and ask students to select the correct expression.

  \item \textbf{Clarify duplicates vs. proper form.} Emphasize that sets do not contain duplicates. If duplicates appear in a written roster (e.g., $\{1,2,3,3\}$), explain that this is not “illegal,” but it is not proper form; the correct form is $\{1,2,3\}$.
  \item \textbf{Avoid Questions That Include Potentially Multiple True Statements.}  The sets $\{1,2,3,3\}$ and $\{2,1,3\}$ are the same set but the first one is not considered in good form. A question oriented toward picking the one in good form is great.  No other questions should entail sets with repeated elements.

  \item \textbf{Keep this lecture introductory.} Focus only on the basics: what sets are, elements/membership, and common representations. Avoid going deep into later topics such as subsets, power sets, unions, intersections, or set differences (these are covered in future lectures).

  \item \textbf{Include at least one infinite-set example.} Provide examples of infinite sets (especially when discussing roster form with ellipses and set-builder form). Briefly contrast \textit{finite} vs.\ \textit{infinite} sets.

  \item \textbf{Apply these boundaries throughout.} Integrate the above guidance naturally wherever relevant in explanations, examples, and exercises.
  \item \textbf{Require correct form.} If a student represents a set incorrectly,such as using a bracket instead of a brace, mark the question incorrect and inform the student of the notation error.  Provide another question to verify the student understands.
  \item \textbf{Examples of bad questions.}
    \begin{itemize}
      \item What is the notation for indicating whether an object is an element of a set?. The problem with this question is the student would have a difficult time showing the symbol. You should provide multiple options for the student to select the appropriate one.
      \end{itemize}
\end{itemize}
